%% =====================================================================
%%  B7-T-13-02 -- what B7 contributes to "Weakness as a Mask"
%%  -------------------------------------------------------------------
%%  LAYER A: APPEND-ONLY. Written once at the entry's write, then left
%%  alone. If the source has more to say later, add a bullet -- do not
%%  rewrite. Material found ONLY here goes in the `unique` slot with
%%  the unique flag set, which makes it undeletable (rule R10).
%% =====================================================================
%% @kind: source
%% @id: B7-T-13-02
%% @book: B7
%% @topic: T-13-02
%% @locator: ch. 4 (pp. 31--34), additive
%% @status: distilled
%% @unique: no
%% @integrated: yes
%% @updated: 2026-09-19

\dossier{B7}{T-13-02}{\bookshort{B7}}{ch. 4 (pp. 31--34), additive}

\begin{distilled}
  \item B7's deception paradox, titled from Congreve --- open truth as the best disguise: the most disarming dodge in their case files is the unsolicited virtue display (the man who answers a theft question with an invitation to inspect his car full of Bibles).
  \item The display works because it pre-credits the speaker's character exactly where an accusation would land; the mask is not weakness here but conspicuous goodness.
  \item B7's rule from the case: the more relevant-seeming and unprovoked the virtue, the more attentive the reader should be to the question it never answered.
\end{distilled}
